\chapter{Referencia técnica}

\section{Mapeo DTMF $\to$ teclado}

\begin{table}[h]
\centering
\begin{tabular}{ccccp{5cm}}
\toprule
\textbf{DTMF} & \textbf{Frecuencias (Hz)} & \textbf{Tecla} & \textbf{Código evdev} & \textbf{Acción en Hugo} \\
\midrule
1 & 697 + 1209 & Q & KEY\_Q & libre / acción 1 \\
2 & 697 + 1336 & $\uparrow$ & KEY\_UP & saltar o subir \\
3 & 697 + 1477 & E & KEY\_E & libre / acción 2 \\
4 & 770 + 1209 & $\leftarrow$ & KEY\_LEFT & izquierda \\
5 & 770 + 1336 & Espacio & KEY\_SPACE & acción principal \\
6 & 770 + 1477 & $\rightarrow$ & KEY\_RIGHT & derecha \\
7 & 852 + 1209 & Z & KEY\_Z & libre / acción 3 \\
8 & 852 + 1336 & $\downarrow$ & KEY\_DOWN & bajar / agacharse \\
9 & 852 + 1477 & X & KEY\_X & libre / acción 4 \\
* & 941 + 1209 & Shift Izq. & KEY\_LEFTSHIFT & modificador \\
0 & 941 + 1336 & Esc & KEY\_ESC & pausa / menú \\
\# & 941 + 1477 & Enter & KEY\_ENTER & confirmar \\
\bottomrule
\end{tabular}
\caption{Mapeo DTMF $\to$ teclado por defecto.}
\end{table}

\section{Archivos clave del sistema}

\begin{table}[h]
\centering
\small
\begin{tabular}{p{0.45\textwidth} p{0.45\textwidth}}
\toprule
\textbf{Archivo} & \textbf{Función} \\
\midrule
\texttt{/etc/asterisk/sip.conf} & Define extensiones SIP 101 y 102 \\
\texttt{/etc/asterisk/extensions.conf} & Dialplan: lógica del 9000 \\
\texttt{/etc/asterisk/manager.conf} & Habilita AMI con usuario \texttt{hugo-bridge} \\
\texttt{/etc/systemd/system/hugo-bridge.service} & Unit del servicio bridge \\
\texttt{/etc/udev/rules.d/99-hugo-uinput.rules} & Permisos sobre \texttt{/dev/uinput} \\
\texttt{/opt/hugo-phone/bridge/hugo-bridge} & Binario compilado del bridge \\
\texttt{/opt/hugo-phone/bridge/config.toml} & Configuración del bridge (hot-reload) \\
\texttt{/opt/hugo-phone/games/hugo/} & Archivos del juego \\
\texttt{\textasciitilde/.config/dosbox/dosbox-staging.conf} & Configuración del emulador \\
\bottomrule
\end{tabular}
\end{table}

\section{Servicios y comandos útiles}

\subsection{Gestión de Asterisk}

\begin{shellcode}
sudo systemctl status asterisk
sudo systemctl restart asterisk
sudo asterisk -rvvv               # consola interactiva
sudo asterisk -rx "reload"        # recargar configs sin reiniciar
\end{shellcode}

Comandos dentro del CLI de Asterisk:

\begin{plaincode}
sip show peers              ; estado de extensiones
sip show registry           ; clientes registrados
sip set debug on/off        ; trazas detalladas
dialplan show hugo-game     ; ver dialplan cargado
manager show settings       ; estado de AMI
manager show users          ; usuarios AMI
manager show connected      ; sesiones AMI activas
core show channels          ; llamadas activas
\end{plaincode}

\subsection{Gestión del bridge}

\begin{shellcode}
sudo systemctl status hugo-bridge
sudo systemctl restart hugo-bridge
sudo journalctl -u hugo-bridge -f      # logs en vivo
sudo journalctl -u hugo-bridge -n 100  # ultimas 100 lineas
\end{shellcode}

\subsection{Recompilar el bridge}

\begin{shellcode}
cd /opt/hugo-phone/bridge
sudo -u beto ~/.cargo/bin/cargo build --release
sudo cp target/release/hugo-bridge ./hugo-bridge
sudo systemctl restart hugo-bridge
\end{shellcode}

\subsection{Inspección del teclado virtual}

\begin{shellcode}
ls /dev/input/by-id/                   # buscar hugo-phone-...
sudo evtest                            # seleccionar el dispositivo
                                       # ver eventos al marcar DTMF
\end{shellcode}

\subsection{Verificación del HT801}

\begin{shellcode}
ping 192.168.1.11                      # alcanzabilidad IP
nmap -p 5060 192.168.1.11              # puerto SIP abierto
\end{shellcode}

\section{Estructura del archivo \texttt{config.toml}}

\begin{tomlcode}
[ami]
host = "127.0.0.1"
port = 5038
username = "hugo-bridge"
secret = "PASSWORD-REAL"

[bridge]
press_duration_ms = 80
event_name = "HugoKey"
key_field = "Key"

[keymap]
"2" = "Up"
"4" = "Left"
"6" = "Right"
"8" = "Down"
"5" = "Space"
"0" = "Esc"
"1" = "Q"
"3" = "E"
"7" = "Z"
"9" = "X"
"*" = "LeftShift"
"#" = "Enter"
\end{tomlcode}

\section{Nombres de teclas soportados}

El parser de \texttt{keyboard.rs} reconoce:

\begin{itemize}
  \item \textbf{Direccionales:} \texttt{Up}, \texttt{Down}, \texttt{Left}, \texttt{Right}
  \item \textbf{Especiales:} \texttt{Space}, \texttt{Enter}, \texttt{Esc} (o \texttt{Escape}), \texttt{Tab}, \texttt{Backspace}
  \item \textbf{Modificadores:} \texttt{LeftShift}, \texttt{RightShift}, \texttt{LeftCtrl}, \texttt{RightCtrl}, \texttt{LeftAlt}, \texttt{RightAlt}
  \item \textbf{Letras:} \texttt{A} a \texttt{Z}
  \item \textbf{Función:} \texttt{F1} a \texttt{F12}
\end{itemize}

Si se asigna un nombre no reconocido, el bridge registra un warning en los logs pero sigue funcionando.

\section{Anatomía del protocolo AMI}

\subsection{Banner inicial}

\begin{plaincode}
Asterisk Call Manager/9.0.0
\end{plaincode}

\subsection{Login}

\begin{plaincode}
Action: Login
Username: hugo-bridge
Secret: PASSWORD-REAL
Events: user,call
(linea en blanco)
\end{plaincode}

\subsection{Respuesta exitosa de login}

\begin{plaincode}
Response: Success
Message: Authentication accepted
(linea en blanco)
\end{plaincode}

\subsection{UserEvent producido por el dialplan}

\begin{plaincode}
Event: UserEvent
Privilege: user,all
UserEvent: HugoKey
Channel: PJSIP/101-00000003
Caller: 101
Key: 2
Uniqueid: 1715641234.3
(linea en blanco)
\end{plaincode}

\section{Glosario}

\begin{description}
  \item[AMI] \emph{Asterisk Manager Interface}. Protocolo TCP de texto plano para controlar Asterisk desde clientes externos.
  \item[ATA] \emph{Analog Telephone Adapter}. Dispositivo que convierte una línea analógica a SIP.
  \item[DOSBox] Emulador del sistema operativo MS-DOS y del hardware del PC compatible. DOSBox-Staging es un fork modernizado.
  \item[DTMF] \emph{Dual-Tone Multi-Frequency}. Sistema de señalización por pares de tonos, usado por los teclados telefónicos.
  \item[evdev] Subsistema del kernel Linux que expone dispositivos de entrada a \emph{user space}.
  \item[FXO] \emph{Foreign eXchange Office}. Puerto que se conecta a una central como si fuera un teléfono.
  \item[FXS] \emph{Foreign eXchange Subscriber}. Puerto que entrega línea como si fuera la compañía telefónica.
  \item[KMSDRM] \emph{Kernel Mode Setting / Direct Rendering Manager}. Capa del kernel Linux que permite dibujar directamente a la pantalla sin servidor X.
  \item[PBX] \emph{Private Branch Exchange}. Centralita telefónica privada.
  \item[RFC 2833] Estándar para transportar tonos DTMF dentro de paquetes RTP, fuera del audio principal.
  \item[RTP] \emph{Real-time Transport Protocol}. Protocolo para transmitir audio y video en tiempo real, sobre UDP.
  \item[SIP] \emph{Session Initiation Protocol}. Protocolo de señalización para iniciar, mantener y terminar sesiones de voz/video.
  \item[uinput] Subsistema del kernel Linux que permite crear dispositivos de entrada virtuales desde \emph{user space}.
\end{description}

\section{Referencias externas}

\begin{itemize}
  \item Documentación oficial de Asterisk: \texttt{wiki.asterisk.org}
  \item DOSBox-Staging: \texttt{dosbox-staging.github.io}
  \item Crate \texttt{uinput} de Rust: \texttt{docs.rs/uinput}
  \item Tokio: \texttt{tokio.rs}
  \item Grandstream HT801 manual: \texttt{grandstream.com/products}
  \item Raspberry Pi documentation: \texttt{raspberrypi.com/documentation}
  \item Hugo en MyAbandonware: \texttt{myabandonware.com/game/hugo-eya}
\end{itemize}

\section{Créditos}

Hugo, el personaje y la franquicia, son propiedad de \textbf{ITE Media ApS} (originalmente \emph{Interactive Television Entertainment}), Dinamarca. Este proyecto es un homenaje técnico sin fines de lucro y sin afiliación oficial con la compañía.

El sistema completo, su código y su documentación fueron desarrollados como proyecto personal por \textbf{Beto} en colaboración con Claude (Anthropic) durante mayo de 2026.

\vspace{1cm}

\begin{center}
\itshape\small
``Hugo, ¡cuidado con la roca!''\\
--- Generaciones de niños chilenos en los noventa
\end{center}
